\documentclass[12pt]{article}
\usepackage[T1]{fontenc}
\usepackage[utf8]{inputenc}
\usepackage[brazil]{babel}
\usepackage{graphicx}
\usepackage{geometry}
\usepackage[export]{adjustbox} 
\usepackage{amsmath, amssymb}
\usepackage{fancyhdr}

% Configurações de layout
\geometry{a4paper, margin=1.5cm, top=2.5cm, bottom=2cm}
\pagestyle{fancy}
\fancyhf{}
\lhead{\textbf{Gabarito Oficial}}
\rhead{Versão: << versao_letra >>}
\cfoot{Página \thepage}

\begin{document}

% --- 1. CABEÇALHO COM LOGO (EXATAMENTE O SEU ORIGINAL) ---
\noindent
\begin{minipage}{0.2\textwidth}
    \includegraphics[width=\textwidth, height=2.5cm, keepaspectratio]{<< logo_path >>}
\end{minipage}
\begin{minipage}{0.75\textwidth}
    \centering
    \textbf{\Large << instituicao >>} \\
    \textbf{Disciplina:} << disciplina_nome >> \\
    \textbf{Profra:} << professor_nome >>

\end{minipage}

\vspace{0.4cm}
\hrule
\begin{center}
    \textbf{\Large GABARITO E RESOLUÇÃO ESPERADA} \\
    \vspace{0.3cm}
    {\Huge \textbf{<< titulo_documento >>}} % Destaque para o Gabarito também
\end{center}
\hrule

\vspace{0.8cm}

% --- 2. LOOP DE QUESTÕES INTELIGENTE ---
<% for q in questoes %>
\noindent \textbf{Questão << loop.index >>} (<< q.pontos >> pts) - \textit{<< q.tipo >>}: 

\vspace{0.2cm}

% Se for Múltipla Escolha ou VF, mostra as alternativas marcando a certa
<% if q.tipo == "Múltipla Escolha" or q.tipo == "Verdadeiro ou Falso" %>
    \begin{itemize}
    <% for alt in q.alternativas %>
        \item [<% if alt.correta %>$\blacksquare$<% else %>$\square$<% endif %>] 
        <% if alt.correta %>\textbf{<< alt.texto >>}<% else %><< alt.texto >><% endif %>
    <% endfor %>
    \end{itemize}
<% endif %>

% Texto da Resolução/Gabarito Discursivo
<% if q.resposta_esperada %>
    \noindent \textbf{Resolução:} \\
    << q.resposta_esperada >>
<% endif %>

% Imagem da Resolução (se existir)
<% if q.gabarito_imagem %>
\begin{center}
    \vspace{0.3cm}
    \textbf{\small Resolução Visual/Cálculo:} \\
    \vspace{0.2cm}
    \includegraphics[max width=0.85\textwidth, max height=10cm, keepaspectratio]{<< q.gabarito_imagem >>}
\end{center}
<% endif %>

\vspace{0.5cm}
\hrule
\vspace{0.5cm}
<% endfor %>

\end{document}